\subsection{Working with Rational Exponents} 
When we evaluate a rational power, we apply the definition. In this case, we will usually use $x^{p/q}=\left(\sqrt[q]{x}\right)^p$.
\begin{example}
Evaluate each of the following, if it is defined.
\begin{lpic}{}{9}
\foreach \y/\exp in {1/{27^{2/3}},
										3/{4^{-3/2}},
										5/{(-32)^{4/5}},
										7/{(-8)^{-1/3}},
										9/{(-81)^{3/4}}}
	\regtextright{3.5}{-\y}{$\exp={}$};
\end{lpic}
\end{example}
\noi We can also do this with variable expressions. In this case, we will usually use $x^{p/q}=\sqrt[q]{x^p}$.
\begin{example}
Rewrite each expression using radicals and positive integer exponents.
\begin{lpic}{}{5}
\foreach \y/\exp in {1/{x^{2/3}},3/{x^{-3/2}},5/{x^{-3/5}}}
	\regtextright{3.5}{-\y}{$\exp={}$};
\end{lpic}
\end{example}
We can apply the usual rules of exponents to rational exponents.
\begin{example}
Simplify each of the following.
\begin{lpic}{}{6}
\foreach \y/\exp in {1/{\left(x^{2/3}y^{3/4}\right)^{3/5}},
										3/{x^{1/3}x^{1/5}x^{1/2}}}
	\regtextright{5.5}{-\y}{$\exp={}$};
	\fractextright{5.5}{-5}{$\left(\dfrac{x^{1/2}y^{2/3}}{x^{1/3}y^{3/4}}\right)^{-12/5}={}$};
\end{lpic}
\end{example}